\section{Operaciones con cardinales}

Dados cardinales $n, m$ (no necesariamente finitos) y conjuntos \emph{disjuntos} $X, Y$ con $\card(X) = n$ y $\card(Y) = m$, se definen:
\begin{enumerate}[label=(\arabic*)]
  \item \textbf{Suma:} $n + m = \card(X \cup Y)$.
  \item \textbf{Producto:} $n \cdot m = \card(X \times Y)$.
  \item \textbf{Potencia:} $n^m = \card\bigl(X^Y\bigr)$, donde $X^Y = \{ f\colon Y \to X \}$.
\end{enumerate}

\begin{note}
  Suponer que $X \cap Y = \varnothing$ no es restrictivo: siempre se pueden reemplazar $X$ e $Y$ por $X \times \{1\} \sim X$ e $Y \times \{2\} \sim Y$, que son disjuntos.
\end{note}

\begin{note}
  Las definiciones no dependen de los conjuntos elegidos. Si $X \sim \tilde{X}$ e $Y \sim \tilde{Y}$, con $X \cap Y = \varnothing = \tilde{X} \cap \tilde{Y}$, y $f\colon X \to \tilde{X}$, $g\colon Y \to \tilde{Y}$ son biyecciones, entonces $h\colon X \cup Y \to \tilde{X} \cup \tilde{Y}$ dada por
  \[
    h(z) =
    \begin{cases}
      f(z) & \text{si } z \in X, \\
      g(z) & \text{si } z \in Y,
    \end{cases}
  \]
  es biyectiva, así que $\card(X \cup Y) = \card(\tilde{X} \cup \tilde{Y})$. Para el producto y la potencia se procede de manera similar.
\end{note}

En lo que sigue, sean $n = \card(X)$, $m = \card(Y)$ y $p = \card(Z)$, con $X, Y, Z$ disjuntos dos a dos.

\begin{prop}
  La suma cumple:
  \begin{enumerate}[label=(\arabic*)]
    \item Conmutatividad: $n + m = m + n$, pues $X \cup Y = Y \cup X$.
    \item Asociatividad: $(n + m) + p = n + (m + p)$.
    \item Neutro: $0 + n = n$, pues $\varnothing \cup X = X$.
  \end{enumerate}
\end{prop}

\begin{prop}
  El producto cumple:
  \begin{enumerate}[label=(\arabic*)]
    \item Conmutatividad: $n \cdot m = m \cdot n$, pues $X \times Y \sim Y \times X$.
    \item $0 \cdot n = 0$, pues $\varnothing \times X = \varnothing$.
    \item $1 \cdot n = n$, pues $\{1\} \times X \sim X$ (la función $(1, x) \mapsto x$ es biyectiva).
    \item Distributividad: $n \cdot (m + p) = n \cdot m + n \cdot p$, pues $X \times (Y \cup Z) = (X \times Y) \cup (X \times Z)$.
  \end{enumerate}
\end{prop}

\begin{note}
  No vale la ley de cancelación: $\aleph_0 + 1 = \aleph_0 + 2$ y $\aleph_0 \cdot 1 = \aleph_0 \cdot 2$, pero $1 \neq 2$.
\end{note}

\begin{eg}
  Sean $\aleph_0 = \card(\N)$ y $c = \card(\R)$. Entonces:
  \begin{enumerate}[label=(\arabic*)]
    \item $\aleph_0 + \aleph_0 = \aleph_0$ y $\aleph_0 \cdot \aleph_0 = \aleph_0$.
    \item $c + c = c$ y $c \cdot c = c$.
  \end{enumerate}
\end{eg}

\begin{proof}
  (1) Los pares $\{2n : n \in \N\}$ y los impares $\{2n - 1 : n \in \N\}$ son numerables, disjuntos y su unión es $\N$; luego $\aleph_0 + \aleph_0 = \aleph_0$. Además $\N \times \N \sim \N$, de modo que $\aleph_0 \cdot \aleph_0 = \aleph_0$.

  (2) Veamos primero que todo intervalo no degenerado tiene cardinal $c$.
  \begin{itemize}
    \item $\arctan\colon \R \to (-\frac{\pi}{2}, \frac{\pi}{2})$ es una biyección, y la función afín $x \mapsto a + (b - a)x$ es una biyección entre $(0, 1)$ y $(a, b)$. Luego $(a, b) \sim \R$ para todo $a < b$.
    \item Si $I$ es un intervalo con extremos $a < b$, entonces $(a, b) \subset I \subset \R$, y por el teorema de Cantor--Schröder--Bernstein, $\card(I) = c$. En particular $(0, 1) \sim [0, 1] \sim \R$.
  \end{itemize}
  Como $(0, 1) = (0, \tfrac{1}{2}] \cup (\tfrac{1}{2}, 1)$ es una unión disjunta de dos intervalos, $c + c = \card(0, 1) = c$.

  Para $c \cdot c = c$ basta ver que $(0, 1) \times (0, 1) \sim (0, 1)$. La función $x \mapsto (x, \tfrac{1}{2})$ es inyectiva de $(0, 1)$ en $(0, 1) \times (0, 1)$. Recíprocamente, escribimos cada $x, y \in (0, 1)$ con su desarrollo decimal infinito (el que no termina en infinitos $9$) $x = 0, x_1 x_2 x_3 \dots$ e $y = 0, y_1 y_2 y_3 \dots$ y definimos $(x, y) \mapsto 0, x_1 y_1 x_2 y_2 \dots$. Esta función es inyectiva, pues el desarrollo obtenido tampoco termina en infinitos $9$, y de él se recuperan $x$ e $y$. Por el teorema de Cantor--Schröder--Bernstein, $(0, 1) \times (0, 1) \sim (0, 1)$.
\end{proof}

\begin{prop}
  $n^m \cdot n^p = n^{m + p}$.
\end{prop}

\begin{proof}
  Hay que ver que $X^Y \times X^Z \sim X^{Y \cup Z}$. La función que a un par $(f, g)$ con $f\colon Y \to X$ y $g\colon Z \to X$ le asigna la función $h\colon Y \cup Z \to X$ dada por $h|_Y = f$ y $h|_Z = g$ (bien definida porque $Y \cap Z = \varnothing$) es biyectiva, con inversa $h \mapsto (h|_Y, h|_Z)$.
\end{proof}

\begin{prop}
  ${(n^m)}^p = n^{m p}$.
\end{prop}

\begin{proof}
  Hay que ver que ${(X^Y)}^Z \sim X^{Y \times Z}$. A cada $F\colon Z \to X^Y$ le asignamos la función $G\colon Y \times Z \to X$ dada por $G(y, z) = F(z)(y)$. Esta correspondencia es biyectiva: su inversa asigna a $G$ la función $F$ dada por $F(z) = G(\cdot, z)$.
\end{proof}

\begin{theorem}
  Sea $n = \aleph_0$ o $n = c$ y sea $m$ un cardinal con $2 \leq m \leq 2^n$. Entonces $m^n = 2^n$.
\end{theorem}

\begin{proof}
  Si $\card(X) \leq \card(Y)$ y $i\colon X \to Y$ es inyectiva, entonces $\varphi \mapsto i \circ \varphi$ es una función inyectiva de $X^Z$ en $Y^Z$, de modo que $\card(X)^{\card(Z)} \leq \card(Y)^{\card(Z)}$. Usando esto,
  \[
    2^n \leq m^n \leq {(2^n)}^n = 2^{n \cdot n} = 2^n,
  \]
  pues $n \cdot n = n$ tanto para $n = \aleph_0$ como para $n = c$. Por el teorema de Cantor--Schröder--Bernstein, $m^n = 2^n$.
\end{proof}

\begin{note}
  ${\{0, 1\}}^{\N} \sim \mathcal{P}(\N)$, pues a cada $f\colon \N \to \{0, 1\}$ se le asigna el subconjunto $A_f = \{ n \in \N : f(n) = 1 \}$ (su función característica es $f$).
\end{note}

\begin{theorem}
  $\R \sim {\{0, 1\}}^{\N}$, es decir, $2^{\aleph_0} = c$.
\end{theorem}

\begin{proof}
  Usaremos el teorema de Cantor--Schröder--Bernstein con $(0, 1) \sim \R$.

  Todo $x \in (0, 1)$ tiene un único desarrollo binario $x = \sum_{n \geq 1} \frac{x_n}{2^n}$ con $x_n \in \{0, 1\}$ que no termina en infinitos $1$. La función $x \mapsto {(x_n)}_{n \geq 1}$ es inyectiva de $(0, 1)$ en ${\{0, 1\}}^{\N}$.

  Recíprocamente, definimos $g\colon {\{0, 1\}}^{\N} \to [0, 1]$ por $g\bigl({(x_n)}_{n}\bigr) = \sum_{n \geq 1} \frac{x_n}{3^n}$. Veamos que es inyectiva. Si ${(x_n)} \neq {(y_n)}$, sea $k$ el primer índice con $x_k \neq y_k$, digamos $x_k = 1$ e $y_k = 0$. Entonces
  \[
    g(x) - g(y) = \frac{1}{3^k} + \sum_{n > k} \frac{x_n - y_n}{3^n} \geq \frac{1}{3^k} - \sum_{n > k} \frac{1}{3^n} = \frac{1}{3^k} - \frac{1}{2 \cdot 3^k} > 0.
  \]
  Luego $g$ es inyectiva (la imagen de $g$ es el conjunto de Cantor, que se verá más adelante). Por el teorema de Cantor--Schröder--Bernstein, $(0, 1) \sim {\{0, 1\}}^{\N}$.
\end{proof}

\section{Hipótesis del continuo}

La \emph{hipótesis del continuo} afirma que no existe ningún cardinal $\kappa$ con $\aleph_0 < \kappa < c = 2^{\aleph_0}$.

Es independiente de los axiomas usuales de la teoría de conjuntos (ZFC, es decir, Zermelo--Fraenkel más el axioma de elección): Gödel probó en 1940 que no puede refutarse y Cohen probó en 1963 que no puede demostrarse.

\section{Construcción de los números reales}

Consideremos sucesiones de números racionales ${(a_n)}_{n \in \N} \subset \Q$. Decimos que ${(a_n)}$ \emph{tiende a $0$}, y escribimos $a_n \to 0$, si para todo $\e > 0$ existe $n_0 \in \N$ tal que $|a_n| < \e$ para todo $n > n_0$.

\begin{definition}[Sucesión de Cauchy]
  Una sucesión ${(a_n)}_{n \in \N} \subset \Q$ es de \emph{Cauchy} si para todo $\e > 0$ existe $n_0 \in \N$ tal que $|a_n - a_m| < \e$ para todos $n, m > n_0$.
\end{definition}

\begin{theorem}
  Si ${(a_n)} \subset \Q$ es convergente, es decir, existe $q \in \Q$ con $a_n - q \to 0$, entonces es de Cauchy.
\end{theorem}

\begin{proof}
  Sea $\e > 0$. Existe $n_0$ tal que $|a_n - q| < \frac{\e}{2}$ para todo $n > n_0$. Entonces, para $n, m > n_0$,
  \[
    |a_n - a_m| \leq |a_n - q| + |q - a_m| < \frac{\e}{2} + \frac{\e}{2} = \e. \qedhere
  \]
\end{proof}

\begin{theorem}
  Toda sucesión de Cauchy ${(a_n)} \subset \Q$ es acotada.
\end{theorem}

\begin{proof}
  Tomando $\e = 1$, existe $n_0$ tal que $|a_n - a_m| < 1$ para todos $n, m > n_0$. En particular, $|a_n - a_{n_0 + 1}| < 1$ para $n > n_0$, de modo que $|a_n| < |a_{n_0 + 1}| + 1$. Luego, con
  \[
    M = \max\{ |a_1|, \dots, |a_{n_0}|, |a_{n_0 + 1}| + 1 \},
  \]
  se tiene $|a_n| \leq M$ para todo $n \in \N$.
\end{proof}

\begin{definition}
  Sea $\mathscr{C}_{\Q}$ el conjunto de todas las sucesiones de Cauchy de números racionales. Dos sucesiones ${(a_n)}, {(b_n)} \in \mathscr{C}_{\Q}$ son \emph{equivalentes} si $a_n - b_n \to 0$.
\end{definition}

\begin{prop}
  La relación anterior es de equivalencia.
\end{prop}

\begin{proof}
  \begin{itemize}
    \item Reflexiva: $a_n - a_n = 0 \to 0$.
    \item Simétrica: si $a_n - b_n \to 0$, entonces $b_n - a_n = -(a_n - b_n) \to 0$.
    \item Transitiva: si $a_n - b_n \to 0$ y $b_n - c_n \to 0$, dado $\e > 0$ existen $n_0$ y $n_1$ tales que $|a_n - b_n| < \frac{\e}{2}$ para $n > n_0$ y $|b_n - c_n| < \frac{\e}{2}$ para $n > n_1$. Para $n > \max\{n_0, n_1\}$ resulta
          \[
            |a_n - c_n| \leq |a_n - b_n| + |b_n - c_n| < \e. \qedhere
          \]
  \end{itemize}
\end{proof}

\begin{definition}[Números reales]
  El conjunto $\R$ de los números reales es el conjunto de clases de equivalencia $[(a_n)]$ de sucesiones de $\mathscr{C}_{\Q}$ respecto de la relación anterior.
\end{definition}
